\documentclass[11pt,a4paper]{article}
\usepackage[utf8]{inputenc}
\usepackage[T1]{fontenc}
\usepackage[french]{babel}
\usepackage[margin=2cm]{geometry}
\usepackage{graphicx}
\usepackage{hyperref}
\usepackage{enumitem}
\usepackage{xcolor}
\usepackage{titlesec}

\setlist{nosep,leftmargin=1.2em}
\titlespacing*{\section}{0pt}{0.6em}{0.3em}
\titlespacing*{\subsection}{0pt}{0.4em}{0.2em}
\titleformat{\section}{\normalfont\bfseries\large}{\thesection}{0.5em}{}
\titleformat{\subsection}{\normalfont\bfseries}{\thesubsection}{0.5em}{}
\setlength{\parskip}{0.3em}
\setlength{\parindent}{0pt}

\title{\vspace{-1.5cm}Rapport individuel --- It\'eration 2}
\author{Louis-Marie Simonneaux (\texttt{lsx4897}) --- Projet Long SHDL}
\date{27 avril 2026}

\begin{document}
\maketitle
\vspace{-1.5em}

\section*{Contexte et objectif de l'it\'eration}

Le projet vise \`a d\'evelopper un simulateur de circuits logiques d\'ecrits en SHDL, avec interface JavaFX. L'\'equipe est r\'epartie sur six axes (\'editeur, parser, automates, simulateur, IHM, int\'egration). Mon axe est la \textbf{repr\'esentation graphique du circuit} : c\^ot\'e utilisateur, la fen\^etre o\`u l'on bascule les entr\'ees et observe les sorties.

L'it\'eration~1 avait livr\'e un prototype JavaFX autonome (deux rang\'ees de \texttt{ToggleButton} ronds et carr\'es) qui d\'emontrait la ma\^itrise des conteneurs JavaFX et du style CSS, mais qui \textbf{n'\'etait reli\'e \`a aucun moteur de simulation}. Le bilan de sprint avait identifi\'e ce manque comme le point bloquant pour rejoindre le MVP bout-en-bout. \textbf{L'objectif de l'it\'eration~2 \'etait donc de brancher cette fen\^etre sur le simulateur r\'eel} (l'API \texttt{Simulateur/Composant/Lien} d\'evelopp\'ee par l'axe simulateur) afin que les changements d'\'etat des entr\'ees se propagent \`a un vrai circuit logique.

\section*{Activit\'es r\'ealis\'ees}

\begin{itemize}
  \item \textbf{Refactoring de \texttt{FenetreSimulation}} (commit \texttt{62c4bbc}, +81/-64 lignes) : r\'eorganisation des m\'ethodes de construction de l'IHM, extraction de la logique de style CSS dans des constantes nomm\'ees (\texttt{StyleUpCarre}, \texttt{StyleDownRond}, \texttt{StyleUnknown\-Carre}), simplification du couplage entre les rangs de \texttt{ToggleButton} et leurs handlers d'\'ev\'enement.
  \item \textbf{Branchement au simulateur} (commit \texttt{6c9bfe4}, 6 fichiers, +131/-41) : la fen\^etre instancie d\'esormais un \texttt{Simulateur(NB\_ENTREES, NB\_SORTIES)} et y ajoute un circuit de d\'emonstration constitu\'e de 4~portes \texttt{And} c\^abl\'ees aux \texttt{Lien} d'entr\'ee et de sortie via \texttt{sim.ajouter\-Composant(...)}. Les clics sur les carr\'es d'entr\'ee modifient les \'etats~UP/DW/ND, le simulateur les propage et les ronds de sortie sont mis \`a jour en cons\'equence.
  \item \textbf{D\'ebut de la classe \texttt{Module}} (\texttt{tests projet long/simulateur/Module.java}) en tant que \texttt{Composant} composite : \texttt{public class Module extends Composant}, conform\'ement \`a la structure pr\'evue par le simulateur. Premi\`ere coquille fonctionnelle, l'impl\'ementation interne (\texttt{calculer()}) reste \`a \'etoffer.
  \item \textbf{Corrections cibl\'ees} pour fiabiliser la batterie de tests c\^ot\'e simulateur : champs \texttt{this.} expli\-cit\'es dans \texttt{Lien.java} (\texttt{4c46433}), correction d'un retour de \texttt{ErreurIndex} et compl\'etude de \texttt{Composant.java} pour les tests JUnit (\texttt{9ef77d7}).
  \item \textbf{Mise \`a jour des scripts de build} (\texttt{LM/Build.sh}) : prise en compte des sources \texttt{tests projet long/simulateur} pour permettre \`a la fen\^etre de compiler avec ses nouvelles d\'ependances.
\end{itemize}

\section*{R\'esultats}

\begin{itemize}
  \item \textbf{Fen\^etre fonctionnelle de bout en bout} : 9~entr\'ees, 9~sorties, 4~portes AND c\^abl\'ees, propagation v\'erifi\'ee manuellement. C'est le premier livrable o\`u l'IHM ne ment plus sur l'\'etat du circuit.
  \item \textbf{Hi\'erarchie de composants \'etendue} : la classe \texttt{Module extends Composant} ouvre la voie \`a la composition de circuits hi\'erarchiques (un module r\'eutilisable comme un composant ordinaire), conform\'ement \`a ce que demande la grammaire SHDL c\^ot\'e parser.
  \item \textbf{Tests existants pass\'es au vert apr\`es modifications} : les ajustements sur \texttt{ErreurIndex} et \texttt{Composant} ont \'et\'e dict\'es par la batterie JUnit du simulateur, qui sert de garde-fou \`a chaque modification structurelle.
\end{itemize}

\section*{Difficult\'es et r\'eussites}

La principale difficult\'e a \'et\'e d'absorber l'API du simulateur sans en \^etre l'auteur : comprendre la s\'emantique des trois \'etats (UP, DW, ND), savoir quand un \texttt{Lien} doit \^etre cr\'e\'e par le \texttt{Simulateur} plut\^ot qu'instanci\'e directement, quel composant d\'eclenche le calcul. La lecture des tests JUnit s'est r\'ev\'el\'ee plus rentable que celle des Javadocs incompl\`etes.

La r\'eussite principale est le fait que la fen\^etre, qui n'\'etait qu'une maquette en sprint~1, soit devenue le \emph{front-end} d'un v\'eritable simulateur. Le c\^ot\'e \og 4~AND hard-cod\'es \fg{} reste assum\'e : c'est un cas test, pas un livrable final ; l'\'etape suivante consiste \`a alimenter le simulateur depuis le module produit par le parser plut\^ot que depuis ce circuit en dur.

\section*{Limites et travaux pour l'it\'eration suivante}

\begin{itemize}
  \item Le circuit interne de la fen\^etre est encore \og en dur \fg{} (4~AND fixes). L'objectif sprint~3 est de remplacer cette construction par un \texttt{Module} produit par la cha\^ine \texttt{parser \(\to\) conversion}.
  \item La classe \texttt{Module} reste un squelette. Sa s\'emantique d\'efinitive d\'epend de la conver\-gence avec l'axe parser/conversion, qui finalisera son contrat \`a l'it\'eration~3.
  \item Aucun branchement avec la \texttt{FenetrePrincipale} (\'editeur de texte SHDL) : la fen\^etre de simulation reste lanc\'ee s\'epar\'ement.
\end{itemize}

\section*{Manuel utilisateur}

\textbf{Pr\'e-requis} : JDK~21+ (testé sur Temurin~25), JavaFX~17 fourni dans le dossier \texttt{javafx-sdk-17.0.18/}.

\textbf{Compiler et lancer la fen\^etre de simulation} (depuis le dossier \texttt{LM/}) :
\begin{verbatim}
  cd LM
  bash Build.sh                         # javac avec module-path JavaFX
  bash Start.sh                         # java SimulationCarresRonds
\end{verbatim}

\textbf{Utilisation} : cliquer sur les carr\'es de la rang\'ee sup\'erieure pour basculer les entr\'ees ; les ronds de la rang\'ee inf\'erieure refl\`etent les sorties calcul\'ees par le simulateur (UP~=~noir, DW~=~blanc, ND~=~point d'interrogation).

\end{document}
